\scanpage{99}
% Image orientation: upside-down typescript; archive foliation: 98.
\uline{Corollaire} 3.5. Soit $G$ un \deleted{\unsure{sché}} préschéma en groupes diagonalisable sur $S$, et soit $n$ un entier $\ne0$. Alors le sous-groupe ${}_n\unsure{H|G}$ de $G$ noyau de l'homomorphisme $n.\operatorname{id}_G:G\quad G$ est entier sur $S$, et fini sur $S$ si $G$ est de type fini sur $S$.

En effet, si $G=D_S(M)$; alors ${}_{\unsure{n}}G=D_S(M/nM)$ en vertu de 3.1., et on conclut par 2.1. c) c').

4. \uline{Fibrés principaux homogènes sous un groupe diagonalisable}.

Soient $S$ un préschéma, et $G=D_S(M)$ un groupe diagonalisable sur $S$. Nous nous proposons de déterminer les \uline{fibrés principaux homogènes} sur $S$, au sens de la ``topologie fidèlement plate quasi-compacte'', cf exposé IV. Comme $G$ lui-même est fidèlement plat et quasi-compact sur $S$, on sait par loc cit que ce sont exactement les $S$-préschémas\inserted{$P$}, \deleted{admettant}\inserted{formellement principaux} avec groupe d'opérateurs $G$, qui sont fidèlement plats et quasi-compacts sur $S$. Par définition, un $S$-préschéma Rappelons qu'un préschéma $P$ sur $S$, à groupes d'opérateurs $G$, est dit principal homogène si tout point de $S$ admet un voisinage ouvert $U$ et un morphisme fidèlement plat et quasi-compact $S'\quad U$, tel que $P'=P\times_S S'$ soit un fibré à opérateurs isomorphe à $G'=G\times_S S'$ (opérant sur lui-même par translations à droite). Comme $G$ est affine sur $S$, il en résulte par SGA VIII 5.6. que $P$ est nécessairement affine sur $S$. Notons aussi que puisque $G$ est lui-même fidèlement plat et quasi-compact sur $S$, $P$ est principal homogène sous $G$ sss il est ``formellement principal homogène'', et s'il est de plus fidèlement plat et quasi-compact sur $S$ (cf exposé IV).

Rappelons d'autre part (exposé I) que la donnée d'un $S$-préschéma $P$ \uline{affine sur} $S$ à groupe d'opérateurs $G=D_S(M)$ revient à la donnée d'une Algèbre quasi-cohérente \uline{graduée de type} $M$ sur $S$, i.e.
